\documentclass[11pt]{article}

\usepackage[margin=1in]{geometry}
\usepackage{amsmath,amssymb,amsthm}
\usepackage{stmaryrd}      % \llbracket \rrbracket
\usepackage[T1]{fontenc}

\newcommand{\PR}{\mathrm{PR}}
\newcommand{\UO}{\mathrm{UO}}
\newcommand{\evalF}{\mathrm{ev}}
\newcommand{\uoValue}{\mathsf{val}}
\newcommand{\bo}{\bot}
\newcommand{\Sw}{S^{\omega}(\bot)}
\newcommand{\N}{\mathbb{N}}
\newcommand{\dotminus}{\mathbin{\dot{-}}}
\newcommand{\D}{D}
\newcommand{\F}{F}
\newcommand{\hatf}{\hat f}

\newtheorem{proposition}{Proposition}
\newtheorem{corollary}{Corollary}
\newtheorem{definition}{Definition}

\title{Ultimate obstination, constructively\\
{\large a machine-checked proof and its computational content}}
\author{}
\date{}

\begin{document}
\maketitle

\begin{abstract}
This note accompanies an Agda formalisation of Coquand's \emph{Une preuve
directe du Th\'eor\`eme d'Ultime Obstination} (\texttt{min1.pdf}), a direct
constructive proof of Colson's ultimate-obstination theorem. We recall the
statement, sketch the proof at a high level, and explain its computational
pay-off: the Scott-continuous extension of a primitive-recursive term is a
computable object, so one can \emph{evaluate} it, in particular at the
infinite diagonal point $x=\Sw$. We illustrate this with a primitive-recursive
$\min$, which returns $\bot$ on $(\Sw,\Sw)$ --- obstination in one line.
Every statement below is machine-checked with
\texttt{agda --safe --without-K --exact-split}, with no postulates, holes, or
pragmas.
\end{abstract}

\section{The domain and primitive recursion}

Let $\D$ be the domain of \emph{lazy naturals}: its elements are the complete
numbers $S^k(0)$, the incomplete finite ones $S^k(\bot)$ (with $S^0(\bot)=\bot$),
and the single infinite element $\Sw$. The order is generated by
$S^k(\bot)\le S^\omega(\bot)$ and $S^k(\bot)\le S^l(x)$ for $k\le l$; the
elements $S^k(0)$ and $\Sw$ are maximal. Write $\F\subseteq\D$ for the finite
elements.

$\PR$ is the usual syntax of primitive-recursive functions: the constant $0$,
the projections, the successor $S$, and closure under composition and
primitive recursion
\[
  f(0,Y)=g(Y),\qquad f(S\,x,Y)=h(x,f(x,Y),Y).
\]
Each $\PR$ term of arity $n$ denotes a monotone map $\F^n\to\F$, written
$\evalF(p)$; moreover every such map is \emph{stable} (Berry): if $A,B$ have a
common upper bound then $f(A\wedge B)=f(A)\wedge f(B)$. Both facts are proved by
a direct induction.

\section{The property and the theorem}

\begin{definition}[Ultimate obstination]
A monotone $f\colon\F^n\to\F$ satisfies the property at $A\in\D^n$ when there is
a finite $A_0\le A$ such that one of:
\begin{enumerate}
\item there is $m$ with $f(X)=S^m(0)$ for all finite $X\ge A_0$;
\item there is $m$ and a coordinate $i$ with $A(i)$ incomplete finite,
      $A_0(i)=A(i)$, and $f(X)=S^m(\bot)$ for all finite $X$ with $X(i)=A_0(i)$
      and $X[i]\ge A_0[i]$;
\item there is a coordinate $i$ with $A(i)=\Sw$, an integer $k$ with
      $A_0(i)=S^k(\bot)$, and a numeric $\varphi$ (constant or strictly
      increasing from $k$) with $f(X)=S^{\varphi(m)}(\bot)$ for all finite $X$
      with $X(i)=S^m(\bot)$, $k\le m$, $A_0[i]\le X[i]$.
\end{enumerate}
The existential and disjunctive content is read intuitionistically. We write
$\UO\,f\,A$ for a witness, and say $f$ \emph{satisfies ultimate obstination} if
it does so at every point.
\end{definition}

In the formalisation this is a three-constructor datatype over the three cases
(file \texttt{Property.agda}).

\begin{proposition}[Colson; Coquand's note]\label{prop:main}
Every element of $\PR$ satisfies ultimate obstination.
\end{proposition}

In Agda (\texttt{Prop1.agda}):
\[
  \mathtt{prop1} : (p:\PR)\,(A:\mathrm{Tup})\to \mathrm{Wf}\,p\,(\mathrm{length}\,A)
  \to \UO\,(\evalF\,p)\,A,
\]
where $\mathrm{Wf}\,p\,n$ records that $p$ is well-formed of arity $n$ (e.g.\ a
projection $\pi_i$ needs $i<n$; $\mathrm{prec}\,g\,h$ at arity $n=m{+}1$ needs
$g$ of arity $m$ and $h$ of arity $m{+}2$).

\paragraph{Proof sketch.} Induction on $p$.
\begin{itemize}
\item \emph{Constants, projections, successor} satisfy the property outright
      (each is eventually constant or reads a single coordinate).
\item \emph{Composition.} The witness functions $\varphi$ are closed under
      composition, so $g\circ\langle h_1,\dots,h_k\rangle$ is obstinate whenever
      the parts are.
\item \emph{Primitive recursion} is the heart. For a finite first argument one
      recurs on its height. For the infinite first argument $x=\Sw$ one forms
      the Kleene approximation
      \[
        u_0=\bot,\qquad u_{k+1}=\hat h(\Sw,u_k,Y),
      \]
      which is increasing and, once $u_k=u_{k+1}$, constant thereafter. A
      constructive dichotomy decides, at each stage $n$, whether the sequence
      has already stabilised (giving a finite limit, hence case 1 or 2) or
      still satisfies $S^n(\bot)\le u_n$. In the second regime one uses that
      $h$ is obstinate at $(\Sw,\Sw,Y)$ and \emph{stable} to pin the growth to
      a single coordinate and read off a witness $\psi$ (with $\psi$ constant
      or strictly increasing), yielding case 3. Stability --- not the stronger
      Kahn--Plotkin sequentiality --- is exactly what excludes two coordinates
      independently driving $f$ to infinity.
\end{itemize}
The recursion argument is carried out once over an abstract base/step pair; the
induction then feeds it the \emph{arity-guarded} interpretations of the
sub-terms (padded to be total and obstinate) and transports the result back to
the concrete interpreter. $\qed$

\section{Computational content: evaluating $\hatf$}

The point of phrasing obstination intuitionistically is that a witness
\emph{carries the value} of the Scott-continuous extension. Here $m$, $i$,
$\varphi$ and $k$ are the data recorded by the witness of Definition~1: $m$ is
the constant of case~1/2, and in case~3 (the infinite coordinate $i$) $k$ is
the threshold with $A_0(i)=S^k(\bot)$ and $\varphi$ is the accompanying numeric
function, which is \emph{either} constant \emph{or} strictly increasing from
$k$ --- case~3 does not fix which; the two sub-cases are read off separately.
Since $A(i)=\Sw=\sup_{m\ge k}S^m(\bot)$ and $f(X)=S^{\varphi(m)}(\bot)$ on the
witnessed region, continuity gives
$\hatf(A)=\sup_{m\ge k}S^{\varphi(m)}(\bot)$, which is $S^{\varphi(k)}(\bot)$
when $\varphi$ is constant (then $\varphi(m)=\varphi(k)$ for all $m\ge k$) and
$\Sw$ when $\varphi$ is strictly increasing. Thus
\[
  \uoValue(\UO\,f\,A)=
  \begin{cases}
    S^m(0) & \text{case 1},\\
    S^m(\bot) & \text{case 2},\\
    S^{\varphi(k)}(\bot) & \text{case 3, }\varphi\text{ constant}
      \ (\text{a fixed finite value}),\\
    \Sw & \text{case 3, }\varphi\text{ strictly increasing},
  \end{cases}
\]
which is the value $\hatf(A)$ of the extension (file \texttt{Property.agda},
$\mathtt{uoValue}$).

\begin{corollary}[Computability]
For every $\PR$ term $f$ of arity $n$ the extension $\hatf\colon\D^n\to\D$ is
computable: $\hatf(A)=\uoValue(\mathtt{prop1}\,f\,A)$ is a total, effective
function of $A$.
\end{corollary}

Since $\mathtt{prop1}$ and $\uoValue$ are total, postulate-free Agda functions,
their composition \emph{is} the algorithm (file \texttt{Computable.agda}):
\[
  \mathtt{fhat}\ :\ (f:\PR)(A:\mathrm{Tup})\to\mathrm{Wf}\,f\,(\mathrm{length}\,A)\to\D .
\]
In particular one may evaluate at the \emph{infinite diagonal} $x=\Sw$ in every
coordinate; call this $\mathtt{fhat\text{-}diag}\,f\,n$. These are not symbolic
answers: for concrete $f$ they reduce, inside the type-checker, to a concrete
element of $\D$.

\section{$\mathrm{pred}$ and $\min$ at $\Sw$}

Take the standard primitive-recursive terms (file \texttt{PredMin.agda}):
\[
\begin{aligned}
  \mathrm{pred} &= \mathrm{prec}\ 0\ \pi_0
     &&\quad \mathrm{pred}(0)=0,\ \mathrm{pred}(S\,x)=x,\\
  \mathrm{sub} &= \mathrm{prec}\ \pi_0\ (\mathrm{pred}\circ\pi_1)
     &&\quad \mathrm{sub}(y,x)=x\dotminus y,\\
  \min &= \mathrm{sub}\circ\langle\, \mathrm{sub}\circ\langle\pi_1,\pi_0\rangle,\ \pi_0\,\rangle
     &&\quad \min(x,y)=x\dotminus(x\dotminus y),\\
  \mathrm{add} &= \mathrm{prec}\ \pi_0\ (S\circ\pi_1)
     &&\quad \mathrm{add}(0,y)=y,\ \mathrm{add}(S\,x,y)=S(\mathrm{add}(x,y)).
\end{aligned}
\]
On total inputs these are genuinely $\mathrm{pred}$, truncated subtraction,
$\min$, and $\mathrm{add}$ (checked by reduction: $\min(2,3)=2$, $\min(3,1)=1$,
$\mathrm{add}(2,3)=5$, etc.). Evaluating their extensions at $x=\Sw$ gives:
\[
  \widehat{\mathrm{pred}}(x)=\Sw,\qquad
  \widehat{\mathrm{sub}}(x,x)=\bot,\qquad
  \boxed{\ \widehat{\min}(x,x)=\bot\ }\ \ (\text{not }x!),\qquad
  \widehat{\mathrm{add}}(x,x)=\Sw.
\]
Each equality holds \emph{definitionally} in Agda ($\mathtt{refl}$).

The reason is exactly ultimate obstination. Subtraction recurs on its first
argument; at $\Sw$ there is no outermost successor to peel, so the Kleene
sequence never rises above $\bot$ and $\widehat{\mathrm{sub}}(\Sw,x)=\bot$.
Since $\min=x\dotminus(x\dotminus y)$ is built from subtraction, it inherits
this obstination and collapses to $\bot$ on $(\Sw,\Sw)$. No primitive-recursive
algorithm for $\min$ can avoid this: it must fully consume an argument, and on
$\Sw$ that computation diverges. This is Colson's theorem, made concrete and
machine-checked.

Addition is the instructive positive contrast. It, too, recurs on its first
argument, but its step function is the \emph{successor}, so the Kleene sequence
$u_0=\bot,\ u_{k+1}=S(u_k)$ is $u_k=S^k(\bot)$, strictly increasing with
supremum $\Sw$. This is case~3 with $\varphi$ strictly increasing, so
$\widehat{\mathrm{add}}(\Sw,\Sw)=\Sw$: infinity plus infinity is infinity. The
denotation faithfully tracks the algorithm --- $\min$ obstinately diverges to
$\bot$, while $\mathrm{add}$ climbs to the infinite element.

\bigskip
\noindent\textbf{Files.}
\texttt{Prop1.agda} (Proposition~\ref{prop:main}),
\texttt{Property.agda} (the property and $\uoValue$),
\texttt{Computable.agda} (the computability corollary),
\texttt{PredMin.agda} (the $\mathrm{pred}/\min$ example),
and the recursion chain \texttt{USeq*}, \texttt{PrecInf*}, \texttt{PrecBot*},
\texttt{PrecAll}. Source note: \texttt{min1.pdf}.

\end{document}
